\documentclass{article}
\newcommand{\scriptopening}[2]{#2}
\newcommand{\scriptframe}[3]{#3}
\newcommand{\ScriptAppendix}{}
\newcommand{\Click}[1]{}
\begin{document}
\scriptopening{0.4 min}{%
  Broadband rollout offers a simple question with a difficult comparison.
  Connected towns may differ from towns that wait, so the timing rule matters
  as much as the outcome pattern. I will show you the assignment first, then
  the estimates. By the end, you should see what the design can establish,
  what remains descriptive, and which evidence would change that judgment.}
\scriptframe{Rollout timing identifies the effect}{0.5 min}{%
  Start with the comparison. Each treated firm is measured against firms in
  towns that have not yet connected. The schedule was recorded before the
  outcome window, and firms did not choose their town's slot. Those facts make
  the timing useful. They do not end the argument, because a predetermined
  score can still favor places on different paths. That is the diagnostic we
  inspect next, and it determines how strongly I will speak.}
\scriptframe{Effects build over three years}{0.6 min}{%
  Read the horizontal axis as years around connection. The first view isolates
  the pre-period, where we ask whether the comparison groups were already
  moving apart. \Click{2} The next view marks arrival. \Click{3} The final view
  adds the post-period estimates. The pattern grows after arrival rather than
  jumping immediately. That delay is consistent with adoption costs, although
  the figure alone cannot identify a mechanism. The design claim therefore
  rests on both halves of this plot, viewed together rather than selectively.}
\ScriptAppendix
\scriptframe{Robustness}{if asked}{%
  The backup repeats the estimate under alternative clustering and comparison
  groups. The bad news is that no specification can repair selection built
  into assignment. The good news is that the reported association is
  numerically stable, which narrows the set of explanations worth testing.}
\end{document}
